\section{Nondegeneracy of the collision-count component}
\label{sec:count-nondegeneracy}
\begin{corollary}[The count component is uniformly nondegenerate]
\label{cor:positive-count-variance}
Let $e_3=(0,0,1,0)$, the collision-count direction. There is a constant
$d_N>0$ such that
\begin{equation}\label{eq:positive-count-variance}
 e_3^{\mathsf T}D_Re_3\ge d_N\qquad(R\in I).
\end{equation}
In particular the actual normalized variance of $N_{n,R}$ is bounded
below by $d_N/2$ for all sufficiently large $n$, uniformly in $R$.
\end{corollary}
\begin{proof}
First, $T_R$ is mixing for the collision probability. The splitting in
Lemma~\ref{lem:collision-spectral-input} proves mixing for smooth
densities and smooth tests. Approximation in $L^2(\nu)$ and invariance
extend it to arbitrary $L^2$ functions: the approximation error in a
correlation is bounded, uniformly in its lag, by Cauchy--Schwarz.
This implication is used for each fixed $R$, without a rate for the
nonsmooth functions below.

Suppose $e_3^{\mathsf T}D_Re_3=0$. By
Theorem~\ref{thm:gaussian-kernel-coboundary}, there is a real
$b\in L^2(\nu)$ with
\[
 1-c_*^{-1}\eta_R=b-b\circ T_R
 \quad\hbox{almost everywhere}.
\]
Set $q=\exp(2\pi i c_*b)$. Since $\eta_R$ is integer-valued,
\[
 q\circ T_R=e^{-2\pi i c_*}q.
\]
Here $0<c_*<1$, so the multiplier $\lambda=e^{-2\pi i c_*}$ is
not one. Invariance gives $\int q\dd\nu=\lambda\int q\dd\nu$,
hence $\int q\dd\nu=0$. Mixing then gives
$\int\overline q(q\circ T_R^m)\dd\nu\to0$. But this integral is
exactly $\lambda^m$, whose modulus is one. This contradiction proves
strict positivity at every $R$. The continuity of $D_R$ and compactness
of $I$ give the uniform constant in \eqref{eq:positive-count-variance}.
Finally \eqref{eq:actual-covariance-rate} identifies the actual count
variance and gives the asserted lower bound for large $n$.

The argument uses the exact integer-valued section indicator and the
nonintegral section mass. It neither chooses a pointwise representative
at a periodic orbit nor proves positive definiteness in every joint
direction.
\end{proof}

